% 第 4 章 §4.2.3 末 + §4.3 + §4.3.1 + §4.3.2（印刷页 185–192，Weil–Atiyah 定理）
% 底本：工程文件/work/en/ch04.txt 第 1201–1723 行。
% 任务书给的范围是第 1201–1684 行，但第 1684 行正落在定理 4.3.5 证明的中间（p.192 尚有 40 行），
% 为让 §4.3.2 收尾完整、环境配平，本块延伸到第 1723 行（p.192 末，§4.3.3 之前）。
% 编号公式：(4.29)(4.30)(4.31)(4.32)(4.33)(4.34)(4.35)(4.36)，共 8 条，无自造编号。
% 定理类环境：引理 4.2.28、引理 4.2.29、定理 4.2.30、命题 4.3.1、定义 4.3.2、定义 4.3.3、
% 命题 4.3.4、定理 4.3.5、注 4.3.6、引理 4.3.7、引理 4.3.8（共 11 个，编号自动生成）。
% 说明：定理 4.2.17 的证明自上一块跨入本块（第 2–4 步），故用 proof 的可选标题承接。

\begin{proof}[定理 4.2.17 的证明（续）]
\emph{第 2 步}（构造非零映射 $\mathcal{E}\to\mathcal{E}_\infty$）。设 $g_j$ 是复规范变换，满足 $g_j(A)=A_j$。$g_j$ 的全纯性蕴含 $\bar\partial_{A_\infty}g_j+[a_j,g_j]=0$。由与第 1 步完全相同的论证，得到估计 $\|g_j\|_{C^\alpha}\le C\|g_j\|_{L^2}$。现在把 $g_j$ 重新缩放，使 $\|g_j\|_{L^2}=1$。上述 $C^\alpha$ 估计对缩放后的映射依然成立，故由紧性可设存在连续映射 $g_\infty\colon\mathcal{E}\to\mathcal{E}_\infty$，使得在 $C^\alpha$ 中 $g_j\to g_\infty$。由于作了归一化，可知 $g_\infty\ne0$。此外，与第 1 步一样可推出 $g_\infty$ 全纯。最后，由 $g_j$ 与 $\Phi_j$ 的 $C^\alpha$ 收敛性以及 $g_j\Phi=\Phi_jg_j$，得到 $g_\infty\Phi=\Phi_\infty g_\infty$。

\emph{第 3 步}（映射 $g_\infty$ 是同构）。假设不然，令 $\mathcal{S}=\ker g_\infty$、$\mathcal{Q}=\mathcal{E}/\mathcal{S}$。则 $\mathcal{Q}$ 是 $\mathcal{E}_\infty$ 的子层。以 $\mathcal{M}$ 记它的饱和化，并令 $\mathcal{N}=\mathcal{E}_\infty/\mathcal{M}$。由于 $\Phi_\infty g_\infty=g_\infty\Phi$，子丛 $\mathcal{S}$ 是 $\Phi$-不变的；同理 $\mathcal{M}$ 是 $\Phi_\infty$-不变的。此外，由 §4.2.1 中的讨论，我们有
\begin{equation}
\begin{aligned}
\mu(\mathcal{Q})-\mu(\mathcal{E}) &\le \mu(\mathcal{M})-\mu(\mathcal{E}),\\
\mu(\mathcal{E})-\mu(\mathcal{S}) &\le \mu(\mathcal{E})-\mu(\mathcal{N}).
\end{aligned}
\label{eq:4.29}\tag{4.29}
\end{equation}
于是得到如下 Higgs 丛的扩张（见 [18]）：
\begin{equation}
\begin{tikzcd}[row sep=large, column sep=large]
0 \arrow[r] & \mathcal{S} \arrow[r] & \mathcal{E} \arrow[r] \arrow[d, "g_\infty"] & \mathcal{Q} \arrow[r] \arrow[d] & 0\\
0 & \mathcal{N} \arrow[l] & \mathcal{E}_\infty \arrow[l] & \mathcal{M} \arrow[l] & 0 \arrow[l]
\end{tikzcd}
\label{eq:4.30}\tag{4.30}
\end{equation}
对 (4.30) 的下行应用引理 4.2.24、对其上行应用引理 4.2.25，得到
\[
\begin{aligned}
&(\operatorname{rank}\mathcal{M})\bigl(\mu(\mathcal{M})-\mu(\mathcal{E})\bigr)
+(\operatorname{rank}\mathcal{N})\bigl(\mu(\mathcal{E})-\mu(\mathcal{N})\bigr)\\
&\qquad\le J(A_\infty,\Phi_\infty)\le\lim_{j\to\infty}J(A_j,\Phi_j)=\inf J(A,\Phi)\\
&\qquad<(\operatorname{rank}\mathcal{S})\bigl(\mu(\mathcal{E})-\mu(\mathcal{S})\bigr)
+(\operatorname{rank}\mathcal{Q})\bigl(\mu(\mathcal{Q})-\mu(\mathcal{E})\bigr),
\end{aligned}
\]
其中既可以用 $J$ 的下半连续性（lower semicontinuity，见 [18]），也可以用 [14, 推论 2.12 与引理 2.17] 中的论证。由于 $\operatorname{rank}\mathcal{M}=\operatorname{rank}\mathcal{Q}$ 且 $\operatorname{rank}\mathcal{S}=\operatorname{rank}\mathcal{N}$，这与 (4.29) 矛盾。

\emph{第 4 步}（Hitchin 方程的解）。最后，我断言 Higgs 丛 $(A_\infty,\Phi_\infty)$ 是 (4.16) 的解。事实上，由 (4.19) 之后的注记，只需证明 $d_{A_\infty}f_{(A_\infty,\Phi_\infty)}=0$ 与 $[f_{(A_\infty,\Phi_\infty)},\Phi_\infty]=0$。后一式成立，因为按假设有 $[f_{(A_j,\Phi_j)},\Phi_j]\to0$（在 $L^2$ 中），而 $f_{(A_j,\Phi_j)}$（相应地 $\Phi_j$）在 $L^p$（相应地 $C^\alpha$）中弱收敛。至于前一式，取 $B$ 为一个试验形式。则
\[
\begin{aligned}
\bigl\langle d_{A_\infty}f_{(A_\infty,\Phi_\infty)},B\bigr\rangle_{L^2}
&=\bigl\langle f_{(A_\infty,\Phi_\infty)},d^*_{A_\infty}B\bigr\rangle_{L^2}\\
&=\lim_{j\to\infty}\bigl\langle f_{(A_j,\Phi_j)},d^*_{A_j}B\bigr\rangle_{L^2}
+\lim_{j\to\infty}\int_X\operatorname{tr}\bigl\{f_{(A_j,\Phi_j)}[a_j,B^*]\bigr\}\\
&=\lim_{j\to\infty}\bigl\langle d_{A_j}f_{(A_j,\Phi_j)},B\bigr\rangle_{L^2}
+\lim_{j\to\infty}\int_X\operatorname{tr}\bigl\{f_{(A_j,\Phi_j)}[a_j,B^*]\bigr\}.
\end{aligned}
\]
第一项消失，因为 $\|d_{A_j}f_{(A_j,\Phi_j)}\|_{L^2}\to0$；第二项也消失，因为 $f_j$ 有界而 $a_j\to0$（在 $C^\alpha$ 中）。由 $B$ 任意，得 $d_{A_\infty}f_{(A_\infty,\Phi_\infty)}=0$，这就完成了证明。
\end{proof}

\begin{proof}[定理 4.2.15 的证明]
设 $[A_j,\Phi_j]$ 是一列多稳定 Higgs 丛，且 $h[A_j,\Phi_j]$ 有界。由定理 4.2.17，可以假设 $(A_j,\Phi_j)$ 满足 (4.16)。由于 $h[A_j,\Phi_j]$ 有界，$\Phi_j$ 的逐点谱一致有界。因此命题 4.2.21 给出 $|\Phi_j|$ 的一致上界。再次利用 (4.16)，得到 $|F_{A_j}|$ 的一致上界。现在可以用 Uhlenbeck 紧性抽出一个收敛子列，它同样满足 (4.16)，正如上面存在性定理的证明中那样。
\end{proof}

\paragraph{Yang--Mills--Higgs 流}
对一对 $(A,\Phi)$，我们用如下方程定义\emph{Yang--Mills--Higgs 流}（Yang--Mills--Higgs flow）：
\begin{equation}
\begin{aligned}
\frac{\partial A}{\partial t}&=-d_A^*\bigl(F_A+[\Phi,\Phi^*]\bigr),\\
\frac{\partial\Phi}{\partial t}&=\bigl[\Phi,\sqrt{-1}\Lambda\bigl(F_A+[\Phi,\Phi^*]\bigr)\bigr].
\end{aligned}
\label{eq:4.31}\tag{4.31}
\end{equation}
上面我们只考虑初值条件中 $\Phi$ 是 $d''_A$-全纯的情形。注意此时这一全纯性条件沿 (4.31) 的解保持不变。事实上，如同 Donaldson [19] 中那样，该流与复规范轨道相切，并且对一切 $0\le t<+\infty$ 存在。这些流方程可以看作 YMH 泛函的 $L^2$-梯度流（gradient flow）；它们推广了 Yang--Mills 流方程。关于这一点更多内容参见 [39, 69] 及其中的参考文献。这里我们只限于讨论几条关键性质；特别是，我们要为上一节中的假设提供依据。

同 (4.17) 那样，令 $f_{(A,\Phi)}=\sqrt{-1}\Lambda(F_A+[\Phi,\Phi^*])$。

\begin{lemma}
对一切 $t\ge0$，
\[
\frac{d}{dt}\operatorname{YMH}(A,\Phi)
=-2\|d_Af_{(A,\Phi)}\|^2_{L^2}-4\|[\Phi,f_{(A,\Phi)}]\|^2_{L^2}.
\]
\end{lemma}

\begin{proof}
我们有
\[
\frac{d}{dt}\operatorname{YMH}(A,\Phi)=2\int_X\operatorname{tr}\bigl(f_{(A,\Phi)}\dot f_{(A,\Phi)}\bigr)\omega.
\]
现在用点表示对时间的导数，
\[
\begin{aligned}
\dot f_{(A,\Phi)}
&=\sqrt{-1}\Lambda\bigl(d_A\dot A+[\dot\Phi,\Phi^*]+[\Phi,\dot\Phi^*]\bigr)\\
&=\sqrt{-1}\Lambda\Bigl(-d_Ad_A^*(F_A+[\Phi,\Phi^*])+\bigl[[\Phi,f],\Phi^*\bigr]+\bigl[\Phi,[\Phi,f]^*\bigr]\Bigr)\\
&=-d_A^*d_Af_{(A,\Phi)}+\sqrt{-1}\Lambda\Bigl([\Phi,f_{(A,\Phi)}]\Phi^*+\Phi^*[\Phi,f_{(A,\Phi)}]+\Phi[\Phi,f_{(A,\Phi)}]^*\\
&\qquad\qquad\qquad\qquad\qquad+[\Phi,f_{(A,\Phi)}]^*\Phi\Bigr).
\end{aligned}
\]
取迹得
\begin{equation}
\operatorname{tr}\bigl(f_{(A,\Phi)}\dot f_{(A,\Phi)}\bigr)
=-\operatorname{tr}\bigl(f_{(A,\Phi)}d_A^*d_Af_{(A,\Phi)}\bigr)
-2\sqrt{-1}\Lambda\operatorname{tr}\bigl([\Phi,f_{(A,\Phi)}][\Phi,f_{(A,\Phi)}]^*\bigr),
\label{eq:4.32}\tag{4.32}
\end{equation}
由分部积分即得结论。
\end{proof}

作为引理 4.2.28 的推论，YMH 沿该流递减。此外还有不等式
\[
\int_0^\infty dt\,\Bigl\{2\|d_Af_{(A,\Phi)}\|^2_{L^2}+4\|[\Phi,f_{(A,\Phi)}]\|^2_{L^2}\Bigr\}
\le\operatorname{YMH}(A_0,\Phi_0).
\]
由此可知：若 $(A_j,\Phi_j)$ 是一列使 $\operatorname{YMH}(A_j,\Phi_j)$ 一致有界的序列，则可把它换成另一列 $(\tilde A_j,\tilde\Phi_j)$，使得 $\operatorname{YMH}(\tilde A_j,\tilde\Phi_j)$ 同样一致有界，但 $d_{A_j}f_{(\tilde A_j,\tilde\Phi_j)}$ 与 $[\Phi_j,f_{(\tilde A_j,\tilde\Phi_j)}]$ 在 $L^2$ 中收敛到 $0$。现在来计算
\[
\begin{aligned}
\Delta\bigl|f_{(A,\Phi)}\bigr|^2
&=-d^*d\bigl|f_{(A,\Phi)}\bigr|^2
=*d*d\operatorname{tr}f_{(A,\Phi)}^2\\
&=2*d*\operatorname{tr}\bigl(f_{(A,\Phi)}d_Af_{(A,\Phi)}\bigr)\\
&=2\operatorname{tr}\bigl(df_{(A,\Phi)}\wedge *d_Af_{(A,\Phi)}\bigr)
-2\operatorname{tr}\bigl(f_{(A,\Phi)}d_A^*d_Af_{(A,\Phi)}\bigr)\\
&=2\bigl|df_{(A,\Phi)}\bigr|^2+4\bigl|[\Phi,f_{(A,\Phi)}]\bigr|^2
+\frac{\partial}{\partial t}\bigl|f_{(A,\Phi)}\bigr|^2,
\end{aligned}
\]
这里用到了 (4.32)。我们已经证明了

\begin{lemma}
对一切 $t\ge0$，
\[
\frac{\partial}{\partial t}\bigl|f_{(A,\Phi)}\bigr|^2-\Delta\bigl|f_{(A,\Phi)}\bigr|^2
=-2\bigl|d_Af_{(A,\Phi)}\bigr|^2-4\bigl|[\Phi,f_{(A,\Phi)}]\bigr|^2.
\]
\end{lemma}

特别地，$\bigl|f_{(A,\Phi)}\bigr|$ 是热方程的下解（subsolution），因此 $\sup\bigl|f_{(A,\Phi)}\bigr|$ 不增。事实上，可以用热核（heat kernel）作显式计算，证明当 $t\ge1$ 时（例如），$\sup\bigl|f_{(A_t,\Phi_t)}\bigr|\le C\operatorname{YMH}(A_0,\Phi_0)$，其中 $C$ 是固定常数。特别地，若 $(A_j,\Phi_j)$ 是一列使 $\operatorname{YMH}(A_j,\Phi_j)$ 一致有界的序列，则可把它换成另一列 $(\tilde A_j,\tilde\Phi_j)$，使 $f_{(\tilde A_j,\tilde\Phi_j)}$ 一致有界。

\begin{proof}[引理 4.2.26 的证明]
在 $(A,\Phi)$ 的复规范轨道中取 $J$ 的一个极小化序列 $(A_j,\Phi_j)$。注意此时 $\operatorname{YMH}(A_j,\Phi_j)$ 一致有界。此外，用与上面类似的论证（见 [14]），$J$ 沿 YMH 流也是递减的。于是，把每个 $(A_j,\Phi_j)$ 换成以 $(A_j,\Phi_j)$ 为初值的 YMH 流上的某一点，同样得到 $J$ 的极小化序列。另一方面，由本节的讨论，我们可以选取流上的点，使条件 (ii) 与 (iii) 也得到满足。这就完成了证明。
\end{proof}

以 $\mathcal{B}^{\min}_E$ 记所有满足 Hitchin 方程 (4.17) 的 Higgs 丛构成的集合。YMH 流建立起一个无穷维的、奇异的 Morse 理论问题，其中 $\mathcal{B}^{\min}_E$ 是泛函的极小值，而不属于 $\mathcal{B}^{\min}_E$、但满足 (4.19) 的 Higgs 丛则扮演更高临界点的角色。可以证明，这一 Morse 理论图景其实不止是类比。特别地，我们有下述

\begin{theorem}[{Wilkin, [69]}]
YMH 流给出 $\mathcal{G}_E$-等变的形变收缩（deformation retraction），把 $\mathcal{B}^{ss}_E$ 收缩到 $\mathcal{B}^{\min}_E$ 上。
\end{theorem}

\section{Betti 模空间}

\subsection{表示簇}

固定基点 $p\in X$，并令 $\pi=\pi_1(X,p)$。以 $\operatorname{Hom}(\pi,\mathrm{SL}_n(\C))$ 记从 $\pi$ 到 $\mathrm{SL}_n(\C)$ 的同态构成的集合。它具有仿射代数簇的结构。令
\[
\mathfrak{M}^{(n)}_B=\operatorname{Hom}(\pi,\mathrm{SL}_n(\C))/\!/\mathrm{SL}_n(\C)
\]
为表示簇，其中双斜杠表示对 $\mathrm{SL}_n(\C)$ 的整体共轭作用取不变量理论的商。于是 $\mathfrak{M}^{(n)}_B$ 是不可约仿射簇，其复维数为 $(n^2-1)(2g-2)$。存在满的代数商映射 $\operatorname{Hom}(\pi,\mathrm{SL}_n(\C))\to\mathfrak{M}^{(n)}_B$，而在不可约（或单）表示构成的开集上它是几何商。$\mathfrak{M}^{(n)}_B$ 的点与半单（或约化）表示的共轭类一一对应；并且 $\operatorname{Hom}(\pi,\mathrm{SL}_n(\C))$ 中每个 $\mathrm{SL}_n(\C)$-轨道的闭包都含有半单表示（这些结果见 [48]）。仿照 Simpson [60, 61]，我把 $\mathfrak{M}^{(n)}_B$ 称为秩 $n$ 的 Betti 模空间（Betti moduli space）。

设 $E\to X$ 是秩 $n$ 的平凡复向量丛。$E$ 上的平坦连接 $\nabla$ 如下给出 $\pi$ 的一个表示。回忆我们已固定基点 $p\in X$，也固定 $E_p$ 的一个标架 $\{e_i\}$。对以 $p$ 为基点的每条环路 $\gamma$，标架 $\{e_i\}$ 的平行移动定义了 $\mathrm{GL}_n(\C)$ 中的一个元素。由于连接平坦，这个元素与同伦类中道路的选取无关。如此便定义了一个元素 $\operatorname{hol}(\nabla)\in\operatorname{Hom}(\pi,\mathrm{GL}_n(\C))$。若 $\nabla$ 在 $\det E$ 上诱导平凡连接，则和乐（holonomy）落在 $\mathrm{SL}_n(\C)$ 中，以下我们总作此假设。反之，给定表示 $\rho\colon\pi_1(X,p)\to\mathrm{SL}_n(\C)$，我们通过商 $\mathcal{V}_\rho=\widetilde X\times\C^n/\pi$ 得到一个带平坦连接 $\nabla$ 的全纯丛 $\mathcal{V}_\rho$，其中 $\widetilde X$ 是 $X$ 的万有覆盖，而该商把 $(x,v)$ 与 $(x\gamma,v\rho(\gamma))$ 等同。

以 $\mathcal{C}_E$ 记 $E$ 上连接构成的空间，$\mathcal{C}_E^{\mathrm{flat}}\subset\mathcal{C}_E$ 为其平坦连接。以 $\mathcal{G}^\C_E(p)$ 记在 $p$ 处为恒同的复规范变换构成的空间，它以共轭作用在 $\mathcal{C}_E$ 上（注意：这与 §4.2.2 中 $\mathcal{G}^\C_E$ 在酉连接空间上的作用不同）。

\begin{proposition}
和乐映射给出 $\mathrm{SL}_n(\C)$-等变的同胚
\[
\operatorname{hol}\colon\mathcal{C}_E^{\mathrm{flat}}/\mathcal{G}^\C_E(p)
\xrightarrow{\ \sim\ }\operatorname{Hom}(\pi,\mathrm{SL}_n(\C)).
\]
特别地，$\mathcal{C}_E^{\mathrm{flat}}/\!/\mathcal{G}^\C_E\simeq\mathfrak{M}^{(n)}_B$。
\end{proposition}

\subsection{局部系统与全纯连接}

\begin{definition}
$X$ 上的复 $n$ 维局部系统（local system）是一个 Abel 群层，它局部同构于常值层（constant sheaf）$\underline{\C}^n$。
\end{definition}

这里 $\underline{\C}$ 表示以 $\C$ 为模型的局部常值层。显然，局部系统 $\mathbf{V}$ 是 $\underline{\C}$ 上的模层。

\begin{definition}
设 $\mathcal{V}\to X$ 是全纯丛。$\mathcal{V}$ 上的\emph{全纯连接}（holomorphic connection）是一个 $\C$-线性算子 $\nabla\colon\mathcal{V}\to K\otimes\mathcal{V}$，满足 Leibniz 法则
\begin{equation}
\nabla(fs)=df\otimes s+f\nabla s,
\label{eq:4.33}\tag{4.33}
\end{equation}
其中 $f\in\mathcal{O}$、$s\in\mathcal{V}$ 为局部截面。
\end{definition}

对一个局部系统 $\mathbf{V}$，令 $\mathcal{V}$ 为全纯丛 $\mathcal{V}=\mathcal{O}\otimes_\C\mathbf{V}$。于是 $\mathcal{V}$ 按如下方式继承一个全纯连接：取 $\mathbf{V}$ 的一个局部平行标架 $\{\mathbf{v}_i\}$。$\mathcal{V}$ 的任一局部截面可以唯一地写成 $s=\sum_{i=1}^n f_i\otimes\mathbf{v}_i$，其中 $f_i\in\mathcal{O}$。然后定义 $\nabla s=\sum_{i=1}^n df_i\otimes\mathbf{v}_i$。由于 $\mathbf{V}$ 的转移函数是常数，这个定义与标架的选取无关；并且 $\nabla$ 立即满足 Leibniz 法则。反之，全纯连接在底层的复向量丛上定义了平坦连接，因为在一个局部全纯标架中曲率 $F_\nabla$ 必为 $(2,0)$ 型，而 Riemann 面上没有 $(2,0)$-形式。特别地，由局部平行截面（即 $\nabla s=0$ 的截面）构成的 $\underline{\C}$-子层 $\mathbf{V}\subset\mathcal{V}$ 定义了一个局部系统。这就给出局部系统与带全纯连接的全纯丛之间的范畴等价（categorical equivalence）（见 [16, 定理 2.17]）。

局部系统有单值表示（monodromy representation）$\rho\colon\pi\to\mathrm{GL}_n(\C)$，它由展开局部平行标架得到。反之，给定 $\rho$，可以像上一节那样构造出一个局部系统。我们有时把二者分别记作 $\mathbf{V}_\rho$ 与 $\mathcal{V}_\rho$。为简单起见，在本书中我几乎总假定单值表示落在 $\mathrm{SL}_n(\C)$ 中；换言之，$\det\mathcal{V}_\rho\simeq\mathcal{O}$，且 $\det\mathcal{V}_\rho$ 上的诱导连接是平凡的。

并非每个全纯丛 $\mathcal{V}$ 都容许全纯连接。特别地，这样的连接是平坦的，故由 (4.3)，$\deg\mathcal{V}=0$ 是必要条件。事实上，关于带全纯连接的丛的 Harder--Narasimhan 型，还可以说得更多。

\begin{proposition}[{cf.\ [8, 23]}]
设 $\mathcal{V}$ 是一个带不可约全纯连接的不稳定丛，并设 $\mu_1>\mu_2>\cdots>\mu_\ell$ 是它的 Harder--Narasimhan 型。则对每个 $i=1,\dots,\ell-1$，有 $\mu_i-\mu_{i+1}\le 2g-2$。
\end{proposition}

\begin{proof}
设 $0\subset\mathcal{V}_1\subset\cdots\subset\mathcal{V}_\ell=\mathcal{V}$ 是 $\mathcal{V}$ 的 Harder--Narasimhan 滤过。由于连接不可约，对每个 $i=1,\dots,\ell-1$，$\mathcal{O}$-线性映射
\[
\mathcal{V}_i\xrightarrow{\ \nabla\ }\mathcal{V}/\mathcal{V}_i\otimes K
\]
都非零。令 $j\le i$ 为使 $\mathcal{V}_j\to\mathcal{V}/\mathcal{V}_i\otimes K$ 非零的最小整数。则由正合列
\[
0\longrightarrow\mathcal{V}_{j-1}\longrightarrow\mathcal{V}_j\longrightarrow\mathcal{Q}_j\longrightarrow 0
\]
可知存在非零映射 $\mathcal{Q}_j\to\mathcal{V}/\mathcal{V}_i\otimes K$。固定这样的 $j$，令 $k\ge i$ 为使 $\mathcal{Q}_j\to\mathcal{V}/\mathcal{V}_k\otimes K$ 非零的最大整数。由
\[
0\longrightarrow\mathcal{Q}_{k+1}\longrightarrow\mathcal{V}/\mathcal{V}_k\longrightarrow\mathcal{V}/\mathcal{V}_{k+1}\longrightarrow 0
\]
可知 $\mathcal{Q}_j\to\mathcal{Q}_{k+1}\otimes K$ 非零。由于诸 $\mathcal{Q}_i$ 都是半稳定的，由引理 4.2.8 得
\[
\mu_j=\mu(\mathcal{Q}_j)\le\mu(\mathcal{Q}_{k+1}\otimes K)=\mu_{k+1}+2g-2,
\]
而结论随之成立，因为 $\mu_i-\mu_{i+1}\le\mu_j-\mu_{k+1}$。
\end{proof}

\paragraph{Weil--Atiyah 定理}
现在的目标是证明下述定理。

\begin{theorem}[{Weil [66], Atiyah [1]}]
全纯丛 $\mathcal{V}\to X$ 容许全纯连接，当且仅当 $\mathcal{V}$ 的每个不可分解（indecomposable）因子的次数为零。
\end{theorem}

这里给出的证明遵循 Atiyah。下述构造将会用得上（见 [1, p.\ 193]）。任何全纯丛 $\mathcal{V}\to X$ 都如下给出一个对应物 $D(\mathcal{V})$。首先，作为光滑丛，$D(\mathcal{V})=(\mathcal{V}\otimes K)\oplus\mathcal{V}$。关于这个分裂，用
\[
f(\varphi,s)=(f\varphi+s\otimes df,fs),\qquad f\in\mathcal{O},\ \varphi\in\mathcal{V}\otimes K,\ s\in\mathcal{V}
\]
定义 $\mathcal{O}$-模结构。可以验证，这使 $D(\mathcal{V})$ 具有 $\mathcal{O}$ 上局部自由层的结构。于是有相容的嵌入 $\varphi\mapsto(\varphi,0)$ 与投影 $(\varphi,s)\mapsto s$，它们使 $D(\mathcal{V})$ 成为一个扩张
\begin{equation}
0\longrightarrow\mathcal{V}\otimes K\longrightarrow D(\mathcal{V})\longrightarrow\mathcal{V}\longrightarrow 0.
\label{eq:4.34}\tag{4.34}
\end{equation}
注意 (4.34) 分裂当且仅当 $\mathcal{V}$ 容许全纯连接。事实上，这样的 $\nabla$ 通过 $s\mapsto(\nabla s,s)$ 给出一个分裂；反之，若 (4.34) 分裂，则存在满足 (4.33) 的 $\C$-线性映射 $\mathcal{V}\to\mathcal{V}\otimes K$。

\begin{remark}
该构造关于子丛是函子性的（functorial）。若 $0=\mathcal{V}_0\subset\mathcal{V}_1\subset\cdots\subset\mathcal{V}_\ell=\mathcal{V}$ 是 $\mathcal{V}$ 的全纯子丛滤过，则有滤过 $0=D(\mathcal{V}_0)\subset D(\mathcal{V}_1)\subset\cdots\subset D(\mathcal{V}_\ell)=D(\mathcal{V})$。
\end{remark}

\begin{lemma}
给定全纯丛 $\mathcal{V}\to X$，设
\[
[\beta]\in H^1\bigl(X,(\mathcal{V}\otimes K)\otimes\mathcal{V}^*\bigr)
\simeq H^{1,1}_{\bar\partial}(X,\operatorname{End}V)
\]
是扩张类。则 $[\operatorname{tr}\beta]=-2\pi\sqrt{-1}c_1(\mathcal{V})$。
\end{lemma}
\begin{annote}
底本自引理 4.3.7 起把花体 $\mathcal{V}$ 排作斜体 $V$（如 $\operatorname{End}V$、定理 4.3.5 证明中的 $\operatorname{End}(V)$、$\deg(V_i)$、“在 $V$ 上取埃尔米特度量”“到 $V_i$ 的正交投影”等处），所指仍是同一个全纯丛 $\mathcal{V}$；译文照排。（来源：上下文推断；Q090）
\end{annote}

\begin{proof}
在 $U_i$ 上取 $\mathcal{V}$ 的局部全纯标架 $s^{(i)}$，并以 $\psi_{ij}$ 记转移函数，$s^{(i)}=s^{(j)}\psi_{i,j}$。我们可以用 $s^{(i)}f^{(i)}\mapsto s^{(i)}\otimes df^{(i)}$ 定义 (4.34) 的局部分裂，其中 $f^{(i)}$ 是 $U_i$ 上一列全纯函数。特别地，
\[
f^{(j)}=\psi_{i,j}f^{(i)},\qquad
\partial f^{(j)}=\psi_{i,j}\bigl(\psi^{-1}_{i,j}\partial\psi_{i,j}f^{(i)}+\partial f^{(i)}\bigr).
\]
由于扩张类由 $I$ 在映射
\[
H^0(X,\operatorname{End}\mathcal{V})\longrightarrow H^1(X,\operatorname{End}\mathcal{V}\otimes K)
\]
下的像给出，由上面的局部分裂可知 $[\beta]$ 由上闭链 $[\psi^{-1}_{i,j}d\psi_{i,j}]$ 代表。于是 $[\operatorname{tr}\beta]=[d\log\det\psi]$。另一方面，若 $h$ 是 $\det\mathcal{V}$ 上的埃尔米特度量，则 $h_i|s^{(i)}_1\wedge\cdots\wedge s^{(i)}_n|^2=h_j|s^{(j)}_1\wedge\cdots\wedge s^{(j)}_n|^2$，故 $h_i|\det\psi_{i,j}|^2=h_j$。这意味着 $d\log\det\psi_{i,j}=\partial\log h_j-\partial\log h_i$。由 Dolbeault 同构，$[\beta]$ 由 $[\bar\partial\partial\log h_i]=[F_{(\bar\partial_{\det\mathcal{V}},h)}]=-2\pi\sqrt{-1}c_1(\mathcal{V})$ 代表（见例 4.2.2 与式 (4.3)）。
\end{proof}

\begin{lemma}
若 $\mathcal{V}\to X$ 是不可分解的全纯丛，且 $\varphi\in H^0(X,\operatorname{End}\mathcal{V})$，则存在 $\lambda\in\C$，使得 $\varphi-\lambda I$ 幂零（nilpotent）。
\end{lemma}

\begin{proof}
由于 $\det(\varphi-\lambda I)$ 全纯而 $X$ 是闭的，$\varphi$ 的特征值必为常数。于是不妨设 $\ker\varphi\ne\{0\},\mathcal{V}$，并考虑正合列
\begin{equation}
0\longrightarrow\mathcal{S}=\ker\varphi\longrightarrow\mathcal{V}\longrightarrow\mathcal{Q}=\operatorname{coker}\varphi\longrightarrow 0.
\label{eq:4.35}\tag{4.35}
\end{equation}
记
\[
\bar\partial_E=\begin{pmatrix}\bar\partial_{\mathcal{S}} & \beta\\ 0 & \bar\partial_{\mathcal{Q}}\end{pmatrix},
\qquad
\varphi=\begin{pmatrix}0 & \varphi_1\\ 0 & \varphi_2\end{pmatrix}.
\]
我们要证明 $\varphi_2=0$。首先注意
\[
0=\bar\partial_E\varphi
=\begin{pmatrix}0 & \bar\partial_E\varphi_1+\beta\varphi_2\\ 0 & \bar\partial_{\mathcal{Q}}\varphi_2\end{pmatrix}.
\]
所以 $\varphi_2$ 作为 $\mathcal{Q}$ 的自同态是全纯的。若 $\varphi_2\ne0$，则它是同构。理由同样是：$\varphi_2$ 的特征值为常数，而按假设 $0$ 不是特征值。于是可以把上面矩阵方程中的右上角元素改写为 $\bar\partial_E(\varphi_1\varphi_2^{-1})+\beta=0$。但那样 $\beta$ 的 Dolbeault 类就为零，(4.35) 也就分裂，这与 $\mathcal{V}$ 不可分解的假设矛盾。
\end{proof}

\begin{proof}[定理 4.3.5 的证明]
设 $\mathcal{V}$ 有全纯连接。则由注 4.3.6，$D(\mathcal{V})$ 分裂。此外，由于 $D(\mathcal{V})$ 关于子丛是自然的，对 $\mathcal{V}$ 的每个不可分解因子都有 $D(\mathcal{V}_i)$ 分裂。但由引理 4.3.7，$\deg(V_i)=0$ 对一切 $i$ 成立。反之，设 $\mathcal{V}$ 不可分解且 $\deg(\mathcal{V})=0$。只需证明 $D(\mathcal{V})$ 分裂。由 Serre 对偶（Serre duality），扩张类为
\[
[\beta]\in H^1\bigl(X,\operatorname{End}(V)\otimes K\bigr)
\simeq\bigl(H^0(X,\operatorname{End}(V))\bigr)^*,
\]
而完美配对（perfect pairing）是 $(\beta,\varphi)=\int_X\operatorname{tr}(\beta\varphi)$。由引理 4.3.8，可以把 $\varphi$ 写成 $\varphi=\lambda I+\varphi_0$，其中 $\varphi_0$ 幂零。于是由引理 4.3.7，
\begin{equation}
\begin{aligned}
(\beta,\varphi)&=(\beta,\varphi_0)+\lambda(\beta,I)
=(\beta,\varphi_0)+\lambda\int_X\operatorname{tr}\beta\\
&=(\beta,\varphi_0)-2\pi\sqrt{-1}\lambda\deg(E)=(\beta,\varphi_0).
\end{aligned}
\label{eq:4.36}\tag{4.36}
\end{equation}
\begin{annote}
底本 (4.36) 第二行作 $\deg(E)$，但本证明自始至终研究的是全纯丛 $\mathcal{V}$（$E$ 是 §4.3.1 中另一个丛的名字），当为 $\deg(\mathcal{V})$ 之误；译文照排 $\deg(E)$。（来源：上下文推断；Q089）
\end{annote}
令 $\mathcal{V}_\ell=\mathcal{V}$，并递归地定义 $\mathcal{V}_{i-1}$ 为 $\varphi_0(\mathcal{V}_i)$ 的饱和化。注意 $\mathcal{V}_{i-1}$ 是 $\mathcal{V}_i$ 的真子丛，否则 $\varphi_0$ 的限制将几乎处处是同构。这个过程终将终止。调整 $\ell$，使 $\mathcal{V}_0=\{0\}$、$\mathcal{V}_1\ne\{0\}$。由注 4.3.6，$\beta$ 保持该滤过 $0=\mathcal{V}_0\subset\mathcal{V}_1\subset\cdots\subset\mathcal{V}_\ell=\mathcal{V}$。在 $V$ 上取埃尔米特度量，并以 $\pi_i$ 记到 $V_i$ 的正交投影。注意
\[
I=\sum_{i=1}^{\ell}(\pi_i-\pi_{i-1})
=\sum_{i=1}^{\ell}(\pi_i-\pi_i\pi_{i-1})
=\sum_{i=1}^{\ell}\pi_i(I-\pi_{i-1}),
\]
且 $(I-\pi_i)\beta\pi_i=(I-\pi_{i-1})\varphi\pi_i=0$。于是
\[
\begin{aligned}
\operatorname{tr}(\beta\varphi_0)
&=\operatorname{tr}(\varphi_0\beta)
=\sum_{i=1}^{\ell}\operatorname{tr}\bigl(\varphi_0\beta\pi_i(I-\pi_{i-1})\bigr)\\
&=\sum_{i=1}^{\ell}\operatorname{tr}\bigl((I-\pi_{i-1})\varphi_0\beta\pi_i\bigr)
=\sum_{i=1}^{\ell}\operatorname{tr}\bigl((I-\pi_{i-1})\varphi_0\pi_i\beta\pi_i\bigr)=0.
\end{aligned}
\]
于是 $(\beta,\varphi_0)=0$，再由 (4.36) 即得 $[\beta]=0$。证毕。
\end{proof}
